\documentclass[../../main.tex]{subfiles}
\begin{document}

{\bf [解]}
題意の2接点を $P\left(t,\dfrac{1}{t}\right)$，$Q\left(-s,\dfrac{1}{s}\right)$ $(t,s>0)$ とする．
すると $P,Q$ での接線$\ell_P$, $\ell_Q$は $\left(\dfrac{1}{x}\right)'=-\dfrac{1}{x^2}$ より
\begin{align}
\ell_P: y=-\frac{1}{t^2}x+\frac{2}{t}, \qquad \ell_Q: y=\frac{1}{s^2}x+\frac{2}{s}
\end{align}
となる．
よって2接線の交点を$R$，$\ell_P$ の $x$ 切片 $T$，$\ell_Q$ の $x$ 切片 $U$ とするとこれらの座標は
\begin{align}
R\left(\frac{2ts(s-t)}{t^2+s^2},\frac{2(t+s)}{t^2+s^2}\right), \qquad T(2t,0), \qquad U(-2s,0) \label{eq:1}
\end{align}
である．
よってグラフは以下のようになる．

\begin{figure}[htb]
\begin{center}
\begin{tikzpicture}[scale=2, >=stealth]
  % 媒介変数の設定 (例: t = 1.2, s = 0.8)
  % P(1.2, 0.833), Q(-0.8, 1.25)
  % T(2.4, 0), U(-1.6, 0)
  % R(-0.384, 1.923)

  % 1. 三角形 RTU の塗りつぶし（面積 f の領域）
  \fill[gray!30] (-1.6, 0) -- (2.4, 0) -- (-0.384, 1.923) -- cycle;
  \fill[pattern=north east lines, pattern color=black!60] (-1.6, 0) -- (2.4, 0) -- (-0.384, 1.923) -- cycle;

  % 2. 座標軸
  \draw[->] (-2.5,0) -- (3.0,0) node[right] {$x$};
  \draw[->] (0,-0.5) -- (0,2.8) node[above] {$y$};
  \node[below left] at (0,0) {$O$};

  % 3. 2曲線（第1象限 y = 1/x, 第2象限 y = -1/x）
  % Q(-s, 1/s) を通るため、x < 0 では y = -1/x のグラフを描画
  \draw[thick, domain=0.45:2.6, samples=100] plot (\x, {1/\x});
  \draw[thick, domain=-2.6:-0.45, samples=100] plot (\x, {-1/\x});

  % 4. 接線 l_P と l_Q (U, R, P, T を通る直線)
  % l_P: y = -0.694 x + 1.667 (U から R を通って T へ)
  \draw[thick, color=blue!80!black] (-1.0, 2.36) -- (2.6, -0.138) node[right] {$\ell_P$};
  % l_Q: y = 1.5625 x + 2.5 (U から R を通って P へ)
  \draw[thick, color=red!80!black] (-1.9, -0.468) -- (0.1, 2.656) node[above right] {$\ell_Q$};

  % 5. 垂線 RH の描画
  \draw[dashed] (-0.384, 1.923) -- (-0.384, 0) node[below, font=\small] {$H$};

  % 6. 各点のプロットとラベル
  \fill (1.2, 0.833) circle (1.2pt) node[above right] {$P$};
  \fill (-0.8, 1.25) circle (1.2pt) node[above left] {$Q$};
  \fill (-0.384, 1.923) circle (1.2pt) node[above] {$R$};
  \fill (2.4, 0) circle (1.2pt) node[below] {$T$};
  \fill (-1.6, 0) circle (1.2pt) node[below] {$U$};

\end{tikzpicture}
\end{center}
\caption{点 $P, Q, R, T, U$ の位置関係と三角形 $RTU$（斜線部）}
\end{figure}

題意の三角形の面積 $f$ は図の斜線部の領域である．$R$ から $x$ 軸への垂足 $H$ として\cref{eq:1}より
\begin{align}
f &=\frac{1}{2}|TU||RH| \\
  &=\frac{1}{2}(2t+2s)\cdot \frac{2(t+s)}{t^2+s^2} \\
  &=\frac{2(t+s)^2}{t^2+s^2} \\
  &=2\left(1+\frac{2ts}{t^2+s^2}\right) \label{eq:2}
\end{align}
である．以下この最大値を求める．
AM-GM不等式 から，
\begin{align}
 t^2+s^2 \ge 2ts \quad (\text{等号成立は}t=s) 
\end{align}
だから
\begin{align}
 \dfrac{2st}{t^2+s^2} \le \dfrac{t^2+s^2}{t^2+s^2} = 1
\end{align}
より，\cref{eq:2}に代入して
\begin{align}
 f \le 2(1+1)=4
\end{align}
である．等号成立は$s=t$の時であり，これは実現可能だから求める面積の最大値は$4$である．

\newpage

\end{document}
